%%%%%%%%%%%%%%%%%%%%%%%%%%%%%%%%%%%%%%%%%%%%%%%%%%%%%%%%%%%%%%%%%%%%%%%%%%%%%%%%
% Medium Length Graduate Curriculum Vitae
% LaTeX Template
% Version 1.2 (3/28/15)
%
% This template has been downloaded from:
% http://www.LaTeXTemplates.com
%
% Original author:
% Rensselaer Polytechnic Institute 
% (http://www.rpi.edu/dept/arc/training/latex/resumes/)
%
% Modified by:
% Daniel L Marks <xleafr@gmail.com> 3/28/2015
%
% Important note:
% This template requires the res.cls file to be in the same directory as the
% .tex file. The res.cls file provides the resume style used for structuring the
% document.
%
%%%%%%%%%%%%%%%%%%%%%%%%%%%%%%%%%%%%%%%%%%%%%%%%%%%%%%%%%%%%%%%%%%%%%%%%%%%%%%%%

%-------------------------------------------------------------------------------
%	PACKAGES AND OTHER DOCUMENT CONFIGURATIONS
%-------------------------------------------------------------------------------

%%%%%%%%%%%%%%%%%%%%%%%%%%%%%%%%%%%%%%%%%%%%%%%%%%%%%%%%%%%%%%%%%%%%%%%%%%%%%%%%
% You can have multiple style options the legal options ones are:
%
%   centered:	the name and address are centered at the top of the page 
%				(default)
%
%   line:		the name is the left with a horizontal line then the address to
%				the right
%
%   overlapped:	the section titles overlap the body text (default)
%
%   margin:		the section titles are to the left of the body text
%		
%   11pt:		use 11 point fonts instead of 10 point fonts
%
%   12pt:		use 12 point fonts instead of 10 point fonts
%
%%%%%%%%%%%%%%%%%%%%%%%%%%%%%%%%%%%%%%%%%%%%%%%%%%%%%%%%%%%%%%%%%%%%%%%%%%%%%%%%

\documentclass[margin]{res}  

% Default font is the helvetica postscript font
\usepackage{helvet}

% Increase text height
\textheight=700pt

\begin{document}

%-------------------------------------------------------------------------------
%	NAME AND ADDRESS SECTION
%-------------------------------------------------------------------------------
\name{Victor Delaplaine}

% Note that addresses can be used for other contact information:
% -phone numbers
% -email addresses
% -linked-in profile

\address{
\\Website : web.calpoly.edu/~vdelapla/
\\LinkedIn : www.linkedin.com/in/victor-delaplaine
\\Github : www.github.com/ByVictorrr\\
}
\address{vdelapla@calpoly.edu \\(+650) 740-6691\\}

% Uncomment to add a third address
%\address{Address 3 line 1\\Address 3 line 2\\Address 3 line 3}
%-------------------------------------------------------------------------------

\begin{resume}

%-------------------------------------------------------------------------------
%	EDUCATION SECTION
%-------------------------------------------------------------------------------
\section{OBJECTIVE}
{\sl Junior Student. Computer Engineering Student at Cal Poly. Inquisitive, hard-working and consistent. Looking for internship opportunities at Tesla where I can apply my skills and contribute to real-world projects }

\section{EDUCATION}
\textbf{California Polytechnic State University}, San Luis Obispo\\
{\sl Bachelor of Science (BS)}, Computer Engineering(CPE)\\
Expected June, 2020
\hfill CGPA: 3.9/4.00 

%-------------------------------------------------------------------------------
%	COMPUTER SKILLS SECTION
%-------------------------------------------------------------------------------
\section{TECHNICAL\\SKILLS}

\textbf{Programming: } Shell, C++, Java, Verilog
\\
\textbf{Database and Markdown :} MySQL, MarkDown, LaTex, html
\\
\textbf{Familiar : } Python, Matlab, CSS, Regular expressions
\\
\textbf{General : } Data Structures, Algorithm, Object Oriented Programming,
\\ \tab Digital Design, Unix/Linux, vim 

%-------------------------------------------------------------------------------
% Modify the format of each position
\begin{format}
\title{l}\\
\dates{l}\location{r}\\
\body\\
\end{format}
%-------------------------------------------------------------------------------

\section{EXPERIENCE}

\textbf{San Francisco State : Machine learning lead intern \hfill{June 18 - August 18}\\}
\normalfont{Developing an android application in java - that reads EMG signals and demodulates them into gestures, output ed to the display. I helped with setting up the EC2 AWS server to get a working SQL base for large amounts of data need to store.}
\begin{itemize}
\item \textbf{Link :} github.com/alexanderadavid/Myo-HMI-Android
\end{itemize}

\section{PROJECTS}

% Project 1
\employer{ByVictorrr}
\location{}
\title{\textbf{aProjector - Autofocusing Projector \hfill September 2018}}
\begin{position}
% \begin{itemize}
  The aProjector, was a invention of me and a peer. I had just taken optics at the time so I was interested and did the physics, getting the equations related it to the the system - programming. I typed up the program in C++ that this program is the ROM on the arduino, it controls the motor to turn the stepper motor a given angle, needed to focus the phone on the frame to focus the image produced on the wall.
% \end{itemize}https://www.overleaf.com/project/5c3cda8d86f90560469d48be
\begin{itemize}
\item \textbf{Technology/Tools:} C++ 
\end{itemize}
\begin{itemize}
\item \textbf{Link :} github.com/byvictorrr/aProjector
\end{itemize}
\end{position}



%Project 2
\employer{ByVictorrr}
\location{}
\title{\textbf{ RPI Camera EC2 Backup - Autonomous backup  \hfill December 2018}}
\begin{position}
% \begin{itemize}
This project was intended for my dad, I built him a security for Christmas using a raspberry pi zero. This camera is viewed from its public IP address - port forwarding port 80. I setup a AWS Linux server to backup the recordings every day to be put on the server, and wipe all the recordings on raspberry pi. This is the lack due to the RPI zero. I learned how to use rsync and cronnie how to schedule shell scripts and remotely backup files to another machine.
% \end{itemize}https://www.overleaf.com/project/5c3cda8d86f90560469d48be
\begin{itemize}
\item \textbf{Technology/Tools: cron, AWS Linux server, ssh, rsync, port forwarding}
\end{itemize}
\begin{itemize}
\item \textbf{Link :} github.com/byvictorrr/RPI-Camera-EC2-Backup
\end{itemize}
\end{position}

%Project 3
\employer{ByVictorrr}
\location{}
\title{\textbf{WindowsServerCalpolyStudents - For UNIX users  \hfill January 2019}}
\begin{position}
% \begin{itemize}
For UNIX UNIX users of any type of operator system that is in need of a windows machine. This could be for anyone who doesn't have enough memory to use a virtual box or dual boot. Intended for Cal Poly EE - students that use mac OS or Linux, because all the EE software is mostly functional on Windows.
% \end{itemize}https://www.overleaf.com/project/5c3cda8d86f90560469d48be
\begin{itemize}
\item \textbf{Technology/Tools: gitbash, RPD - protocol, AWS, powershell}
\end{itemize}
\begin{itemize}
\item \textbf{Link :} github.com/byvictorrr/WindowServerCalpolyStudents
\end{itemize}
\end{position}

\section{CERTIFICATION}
\par
\normalfont{}
\\
\normalfont{\textbullet{} \textbf{Computer Science - C++ Certificate of Achievement} by Ca\tilde{n}ada College
\\ 
\textbullet{} \textbf{Engineering - Certificate of Achievement} by \hspace{2pt} Ca\tilde{n}ada \hspace{2pt}College 


\section{RELEVANT\\COURSES}
\par

\normalfont{\textbullet{} Digital Design: github.com/byvictorrr/CPE133 }
\\
\normalfont{\textbullet{} Object Orientated Programming in Java: github.com/byvictorrr/CPE203 }
\\
\normalfont{\textbullet{} Object Orientated Programming in C++: https://github.com/ByVictorrr/Sudoku }
\\
\normalfont{\textbullet{} Data Structures in C++: https://github.com/ByVictorrr/GhostGame}
\\



\section{ADDITIONAL ACTIVITIES}

\normalfont{ \textbullet{} , I like to watch videos on vim scripting}

%-------------------------------------------------------------------------------


\end{resume}
\(\)\end{document}
